\subsection{APPLICATION TO ELEMENTARY AUTOMORPHISMS}

PROPOSITION 9. Let k be a field of characteristic 0 and g a Lie algebra over k. If \(a \in \operatorname{Aut}_{e}(\mathfrak{g})\) , the dimension of the nilspace of a - 1 is greater than or equal to the rank of g.

By the ``Lefschetz principle'' (Algebra, Chap. V, §14, no. 6, Cor. 2 of Th. 5), k is an ascending directed union of subfields \((k_{i})_{i\in\mathrm{I}}\) which admit C as extension field. Let \((e_{\alpha})\) be a basis of g over k and \(x_{1},\ldots,x_{m}\) elements of g such that \(\operatorname{ad}(x_{1}),\ldots,\operatorname{ad}(x_{m})\) are nilpotent and \(a=e^{\operatorname{ad}(x_{1})}\ldots e^{\operatorname{ad}(x_{m})}\) . Let \(c_{\alpha\beta}^{\gamma}\) be the structure constants of g with respect to the basis \((e_{\alpha})\) and \((x_{r}^{\alpha})\) the components of \(x_{r}\) with respect to this basis \((1\leq r\leq m)\) . There exists an index \(j\in I\) such that the \(c_{\alpha\beta}^{\gamma}\) and the \(x_{r}^{\alpha}\) all belong to \(k_{j}\) . Let \(g_{j}=\sum_{\alpha}k_{j}e_{\alpha}\) ; this is a Lie algebra over \(k_{j}\) containing \(x_{1},\ldots,x_{m}\) , and the restriction \(a_{j}\) of a to \(g_{j}\) is an elementary automorphism of \(g_{j}\) . The extension of \(a_{j}\) to \(g_{j}\otimes_{k_{j}}C\) is an elementary automorphism \(a_{j}\otimes1\) of \(g_{j}\otimes C\) . So let \(G_{j}\) be a connected complex Lie group with Lie algebra \(g_{j}\otimes C\) , and s an element of \(G_{j}\) such that \(\operatorname{Ad}(s)=a_{j}\otimes1\) . Prop. 8, applied to the pair \((G_{j},s)\) , shows that the nilspace of \(a_{j}\otimes1-1\) is of dimension n, so

\[
n \geq \operatorname{rk} (\mathfrak {g} _ {j} \otimes \mathbf {C}) = \operatorname{rk} (\mathfrak {g} _ {j}) = \operatorname{rk} (\mathfrak {g}).
\]

But this nilspace has the same dimension as that of \(a_{j} - 1\) and that of \(a - 1\) . Hence the proposition.

